\chapter{Research Design and Problem Formulation}\label{research-design-and-problem-formulation}

\section{Research Strategy}\label{research-strategy}

The thesis follows a design-and-evaluation strategy. It first defines a
verification contract and implements a modular system that realizes it. It then
constructs a reviewed benchmark and compares system configurations through
controlled experiments. Quantitative evaluation measures requirement labels,
evidence-step retrieval, cost, and stability. Qualitative analysis explains
which requirement structures and evidence patterns produce errors. Region
grounding and UI evaluability are evaluated as auxiliary components because
they affect traceability and the visible scope of the task.

The central unit of analysis is a verification item: one textual requirement
paired with one ordered screenshot flow and one reviewed reference decision.
Items from the same flow are not independent because they share screenshots and
application context. Statistical comparisons therefore use a cluster bootstrap
that resamples complete flows rather than individual requirements (Field and
Welsh, 2007). With only 13 flow clusters, intervals are
reported as uncertainty summaries and interpreted cautiously.

The main evaluation is retrospective. The system does not control the website
or request new states; it judges an existing sequence. This choice makes the
visual evidence reproducible but limits conclusions to the recorded path. It
also makes it possible to evaluate failures of screenshot selection separately
from failures of visual interpretation.

\section{Formal Task Definition}\label{formal-task-definition}

Let an ordered screenshot flow be

\[
F = \langle s_1, s_2, \ldots, s_n \rangle,
\]

where each \(s_i\) contains an image, a stable step index, and optional textual
metadata such as OCR or extracted page text. Let \(r\) be a textual
requirement. The verifier implements a mapping

\[
V(r, F) \rightarrow (y, a, E, C, U, q),
\]

where \(y\in Y\) is a requirement-level label, \(a\in A\) is a UI-evaluability
decision, \(E\) is a set of evidence units, \(C\) is an optional set of
claim-level decisions, \(U\) contains uncertainty reasons, and \(q\) is a
rationale. The labels belong to

\[
Y = \{\texttt{FULFILLED},\ \texttt{PARTIALLY\_\allowbreak FULFILLED},\
\texttt{NOT\_FULFILLED},\ \texttt{ABSTAIN}\}.
\]

\[
A = \{\texttt{UI\_VERIFIABLE},\
\texttt{PARTIALLY\_\allowbreak UI\_\allowbreak VERIFIABLE},\
\texttt{NOT\_UI\_VERIFIABLE}\}.
\]

Together, \(y\) and \(a\) distinguish two reasons why a requirement may not
receive \texttt{FULFILLED}: the recorded flow may omit evidence for an otherwise
UI-verifiable requirement, or the requirement may contain obligations that
screenshots cannot resolve in principle.

An evidence unit minimally identifies a screenshot step and may additionally
contain a textual observation and one or more regions, represented by bounding
boxes and region metadata. The ordering of the flow is part of the input. Two
flows containing the same screenshots in a different order do not necessarily
express the same interaction history.

The target is the conclusion supported by the recorded visual evidence. It does
not cover every possible execution of the implementation. The label therefore
depends on the written requirement and the evidence policy supplied to
annotators and models.

\section{UI Evaluability}\label{ui-evaluability}

UI evaluability answers whether visible UI evidence can in principle resolve
the material obligations of a requirement. It is annotated separately from
fulfillment:

\begin{itemize}
\tightlist
\item
  \texttt{UI\_VERIFIABLE} means that the material claims can be assessed through
  rendered interface content or visible state changes.
\item
  \texttt{PARTIALLY\_\allowbreak UI\_\allowbreak VERIFIABLE} means that a material visible core exists but full
  satisfaction also depends on hidden state, persistence, policy, business
  logic, or an external system.
\item
  \texttt{NOT\_UI\_VERIFIABLE} means that the central property has no stable visual
  manifestation or is too abstract for screenshot-based assessment.
\end{itemize}

For the complete operational definitions, uncertainty reasons, and consistency
gates, see Appendix A.

Evaluability is a property of the requirement relative to the observation
modality, not a synonym for evidence availability in one particular flow. A
requirement can be UI-verifiable even when the recorded flow omits the needed
state. Its fulfillment label then depends on the evidence that remains;
\texttt{ABSTAIN} is appropriate only when the flow supports no reliable positive,
partial, or negative decision. Conversely, a partially UI-verifiable
requirement may receive a partial label when its visible part is supported
while a material hidden obligation remains unresolved.

Separating evaluability from fulfillment keeps hidden properties outside the
visible contract while preserving the difference between an unobservable
requirement and an incomplete trajectory. Chapter 6 measures agreement with the
evaluability schema and stratifies verification errors by class.

\section{Verification Label Semantics}\label{verification-label-semantics}

\texttt{FULFILLED} is the strongest claim. It requires visible support for every
material UI-observable obligation, at least one recorded evidence unit, no
visible contradiction, and no unresolved material uncertainty about the
visible behavior. Routine implementation dependencies do not block fulfillment
when the requirement is explicitly satisfied by a visible success proxy. A
cart badge update, for example, can establish the requested visible cart state;
it does not prove database durability.

\texttt{PARTIALLY\_\allowbreak FULFILLED} describes a compound evidence state: at least one
important claim has visible support, while another remains missing, hidden, or
ambiguous. For a requirement covering both a search form and complete correct
results, the form can be supported even when result correctness and
completeness remain unverified.

\texttt{NOT\_FULFILLED} requires visible counter-evidence against a central observable
claim. Missing evidence is insufficient. A requested control that is visibly
absent from the relevant complete state may justify the negative label, whereas
a flow that never visits that state does not.

\texttt{ABSTAIN} is used when the evidence cannot support a reliable positive, partial,
or negative decision. Reasons include textual ambiguity, unclear scope,
unresolved quantifiers, missing before/after states, unverified outcomes, and
nontrivial hidden properties. An abstention is a completed semantic prediction,
not a missing response.

\section{Claims and Deterministic Aggregation}\label{claims-and-deterministic-aggregation}

Long requirements may contain several obligations connected by conjunctions,
conditions, or outcome clauses. The pipeline may decompose a requirement into
claims \(C_r = \{c_1,\ldots,c_m\}\). Claim text is derived from the requirement
before the screenshots are interpreted. It must not encode what a particular
screen happens to show, because doing so would leak the evidence into the
contract.

Each claim receives one of the evidence statuses \texttt{SUPPORTED},
\texttt{CONTRADICTED}, \texttt{MISSING}, \texttt{HIDDEN}, \texttt{AMBIGUOUS}, or \texttt{OUT\_OF\_SCOPE}. Claims may
also be marked as core or supporting obligations. A deterministic aggregator
then maps the collection of statuses and the requirement's UI evaluability to a
requirement label. In simplified form:

\begin{itemize}
\tightlist
\item
  all observable core claims supported and evidence present implies
  \texttt{FULFILLED};
\item
  supported and unresolved important claims imply \texttt{PARTIALLY\_\allowbreak FULFILLED};
\item
  a contradicted observable core claim implies \texttt{NOT\_FULFILLED};
\item
  insufficient, hidden, or ambiguous evidence without meaningful partial
  support implies \texttt{ABSTAIN}.
\end{itemize}

The aggregation rules standardize final-label interpretation but cannot repair
incorrect claim statuses. They encode the additional policy choice that a
negative decision requires visible contradiction. An offline counterfactual
tests the alternative of converting native abstentions to negative labels.

\section{Evidence Contract}\label{evidence-contract}

Evidence is represented at two nested granularities. Step-level evidence
identifies the relevant screenshot or transition. Region-level evidence can
then refine that trace by localizing an observation within a selected
screenshot. In the final multimodal runs, step indices, textual observations,
and optional regions are returned jointly in the same model response, but their
quality is evaluated separately:

\begin{enumerate}
\def\labelenumi{\arabic{enumi}.}
\tightlist
\item
  Did the system identify the correct state or transition?
\item
  Within that state, did it identify a relevant and sufficient visual region?
\end{enumerate}

Step-level traceability is evaluated through overlap between returned and
reviewed evidence steps. The evaluator reports hit@k, recall@k, precision@k,
and a reciprocal first-hit summary. Requirement-level evidence is normalized
into chronological step order, so position-based values describe alignment of
the ordered trace rather than a confidence or lexical-retrieval ranking.
Region-level evidence is evaluated by coordinate validity, applicability,
coverage, relevance, and sufficiency.
Intersection over union is useful when a reference rectangle exists, but it is
not sufficient because several spatially different boxes may all provide valid
evidence.

Some claims are not reducible to one rectangle. A comparison may need two
regions, a transition may need two screenshots, and a screen-wide state may not
have a meaningful minimal box. The region-evaluation protocol therefore permits
\texttt{SINGLE\_REGION}, \texttt{MULTI\_REGION}, \texttt{WHOLE\_\allowbreak SCREEN\_\allowbreak OR\_\allowbreak TRANSITION}, and
\texttt{NO\_\allowbreak VISIBLE\_\allowbreak REGION}. Returning no region is treated as a localization
abstention. It can be appropriate when the claim has no localizable visible
support; otherwise it is a localization failure.

\section{Operationalization of the Research Questions}\label{operationalization-of-the-research-questions}

Matched full-coverage runs on the 258-item benchmark answer RQ1. Accuracy,
macro-F1, class-specific scores, confusion matrices, false fulfillment,
abstention, agreement, cost, and run stability cover different aspects of
fulfillment assessment. The deterministic and hosted open-weight baselines
provide reference points. RQ1 additionally compares matched configurations
across different models to assess model sensitivity.

For RQ2, a controlled two-by-two matrix crosses raw versus gated-decomposed
requirements with complete-flow versus shared lexical top-4 screenshots. Model, prompt
family, label schema, aggregation, batching, and benchmark remain fixed.
Screenshot-step traceability and cost are evaluated alongside labels. The
separate region-grounding analysis does not enter the accuracy comparison.

RQ3 applies a predefined taxonomy to frozen predictions. A screenshot-aware
coding pass assigns one primary outcome and optional pattern tags, while a
deterministic heuristic serves only as a separate consistency check. Counts are
descriptive for the reviewed flows and do not support population-level or
causal claims.

\section{Scope and Non-Claims}\label{scope-and-non-claims}

The task covers visible content, controls, navigation outcomes, short state
transitions, validation, and visible interaction results. It does not establish
hidden backend correctness, security, external delivery, global availability,
long-term persistence, or strict performance properties unless a requirement
concerns only their visible representation. Accordingly, the study provides
neither formal verification nor exhaustive testing; it evaluates what a finite
recorded flow supports.
